\documentclass[../../../include/open-logic-section]{subfiles}

\begin{document}

\olfileid{sth}{spine}{valpha}

\olsection{تعریف مراحل به‌منزلهٔ $V_\alpha$ها}

در~\olref[sth][ordinals][]{chap} خوش‌ترتیب‌ها و اعداد ترتیبی (فون نویمان)
را تعریف کردیم. در این فصل از آن‌ها برای توصیف خودِ سلسله‌مراتب مجموعه‌ها
استفاده خواهیم کرد. برای این منظور به یاد آورید که در
\olref[sth][ordinals][opps]{sec} مفاهیم اعداد ترتیبی جانشین و حدی را تعریف
کردیم. این مفاهیم را در تعریف زیر به کار می‌بریم:

\begin{defn}\ollabel{defValphas}
	\begin{align*}
	V_\emptyset &\defis \emptyset\\
	V_{\ordsucc{\alpha}} &\defis \Pow{V_\alpha} & &
	\text{برای هر عدد ترتیبی }\alpha\\
	V_{\alpha} &\defis \bigcup_{\gamma < \alpha} V_\gamma & &
	\text{اگر }\alpha\text{ عدد ترتیبی حدی باشد}
\end{align*}
\end{defn}
\noindent
این تعریفی با \emph{بازگشت ترامتناهی} روی اعداد ترتیبی خواهد بود. از این
حیث باید آن را با تعریف‌های بازگشتیِ توابع روی اعداد طبیعی مقایسه
کنیم.\footnote{برای نمونه، تعریف‌های جمع، ضرب و توان‌رسانی را در
\olref[sfr][infinite][dedekind]{sec} ببینید.} همانند کار با اعداد طبیعی،
یک حالت پایه و حالت‌های جانشین تعریف می‌کنیم؛ اما هنگام کار با اعداد
ترتیبی باید رفتار حالت‌های \emph{حدی} را نیز توصیف کنیم.

این تعریفِ $V_\alpha$ها نقطهٔ عطف مهمی خواهد بود. به‌طور غیرصوری انگیزهٔ
سلسله‌مراتب مجموعه‌ها را فرایندی دانستیم که مجموعه‌ها را در
\emph{مراحل} تشکیل می‌دهد. در واقع، $V_\alpha$ها همان مراحل‌اند. بااین‌حال،
نکتهٔ مهم آن است که این توصیفی \emph{درونی} از مراحل است. به‌جای آنکه صرفاً
\emph{مدلی} ممکن برای نظریه پیشنهاد کنیم، مراحل را \emph{درون} نظریهٔ
مجموعه‌های خود تعریف خواهیم کرد.

\end{document}
